\section{征服之后，行省是接收而不是终点}

1260--1307年的连续纪年跨越忽必烈即位、1267年营建大都、1271年定国号元、1276年进入临安、1277年在江南设置行省与财赋接收安排、1279年南宋政权终结，以及1294年忽必烈去世后的继承。事件顺序可由《元史》本纪建立，但明初仓促编成的文本有纪年、编次与终结视角的问题；它也不能代替蒙古文、波斯文、高丽材料或江南地方文书。

行省、路府与财赋接收的名目说明国家试图如何组织空间，不等于一处地区从此被同样管理。征服后的官员、军户、匠户、站户、寺院和城市居民在材料中出现的方式不同；“色目”“汉人”等分类更须按具体法律、时期与区域说明，不能被当作稳定的社会本质。江南的交接也不是单向下达，而涉及仓储、漕运、地方精英、道路与反抗的多重协商。

\section{纸钞、站赤与地方的可见性}

《元典章》《通制条格》及行省文书适合追问机构、规范和反复申令；它们不自动证明实际执行。纸钞发行额不等于流通，更不能直接替代物价、赋役或家庭生计；站赤、河工和军户的材料也各有地域与身份边界。讨论元的财政行政，应将条格、官文书、碑刻、契约和考古分别放在它们能证明的位置。

1267年以后大都的营建、1277年江南接收、1281年对日远征受挫、1287年乃颜叛乱和1294年继承，显示中枢始终须在北方、江南、东北与海域之间调度资源。不能把大都或江南的经验当作全帝国的缩影。元朝国家能力最值得追问的，不是它是否拥有一个统一称号，而是不同语言、道路、税役和地方机构如何在不平衡的材料中被看见。
